% ======================================================================
%                  CONCLUSIONES Y RECOMENDACIONES
% ======================================================================
\chapter*{CONCLUSIONES Y RECOMENDACIONES}
\addcontentsline{toc}{chapter}{CONCLUSIONES Y RECOMENDACIONES}\label{cap:conclusiones}

\section*{Conclusiones}

\begin{enumerate}
\item La problemática se delimitó como una limitación de ingeniería: los planes de tiempo fijo o con
adaptación limitada no representan adecuadamente las variaciones direccionales de demanda, mientras que una
arquitectura exclusivamente remota introduce dependencia de la conectividad y un punto único de falla. Esto
justifica un módulo local complementario y no el reemplazo directo del controlador semafórico (OE1, OE2).

\item El diseño conceptual estableció una cadena trazable entre necesidades, requerimientos, funciones y
alternativas. La evaluación técnico-económica seleccionó el concepto cámara + procesamiento en el borde +
ICV + control difuso + nube segura + gabinete modular, principalmente por su autonomía local, su
interpretabilidad y su compatibilidad con una validación reproducible (OE3).

\item El diseño electrónico definió y dimensionó las funciones de sensado, procesamiento, comunicación,
alimentación, protección, respaldo y diagnóstico. El balance de potencia (consumo cercano a 145.7~W,
elevado a 182~W con un factor de seguridad del 25~\%), el dimensionamiento de las fuentes, la autonomía del
respaldo ($\sim$72~min) y la coordinación de protecciones quedaron resueltos a nivel de diseño; su cierre
para fabricación requiere la confirmación con componentes físicos y con el controlador objetivo (OE4).

\item El diseño mecánico incorporó gabinete, distribución interna, ventilación, acceso de mantenimiento y
soporte de cámara. El análisis estructural por elementos finitos arrojó un factor de seguridad mínimo de
5.94 y un desplazamiento máximo de 0.21~mm bajo la carga de viento de diseño, y el análisis térmico estimó
un incremento de 2.5~\textdegree C; la clasificación IP, la resistencia al vandalismo y el comportamiento
térmico en campo deben confirmarse mediante prototipo y ensayos (OE5).

\item El control difuso local se definió como núcleo por su interpretabilidad y su capacidad de operar en el
borde. Las entradas, las funciones de pertenencia, las reglas, la inferencia, la defuzzificación y las
saturaciones permiten revisar el comportamiento antes de emitir una recomendación. Ante una falla del
sensado o del módulo, la arquitectura inhibe la adaptación y conserva el plan seguro del controlador, que
mantiene la autoridad final sobre la actuación (OE6).

\item La nube se delimitó como soporte para telemetría, históricos, alertas, coordinación y distribución
autorizada de parámetros, demostrándose que la decisión local continúa sin conectividad externa. La
ciberseguridad propuesta incorpora autenticación, cifrado, control de acceso por roles, validación de
parámetros, registro de eventos y mínimo privilegio; estas medidas son criterios de diseño y no equivalen a
una auditoría o certificación (OE7, OE8).

\item La validación en SUMO comparó el control de tiempo fijo con el control difuso adaptativo en el
escenario de Lima-Centro, bajo condiciones comunes y con tres semillas, conservando el mismo presupuesto de
verde por ciclo. El control difuso redujo la demora media de 504.3 a 477.1~s/veh ($-5.4$~\%), la saturación
de cola en $3.6$~\%, el índice de congestión de 0.5055 a 0.4954 ($-2.0$~\%), e incrementó la velocidad media
en $1.4$~\% y el \textit{throughput} en $1.8$~\%, manteniendo en todo momento los límites duros de verde
$[10,120]$~s. Las mejoras son modestas pero consistentes en promedio; en una de las tres semillas el control
difuso no superó al plan fijo, lo que evidencia que el beneficio depende de la demanda y no está garantizado
en toda condición. Estos resultados validan el núcleo de control y, a la vez, acotan honestamente su alcance
(OE9).

\item El análisis económico preliminar separó CAPEX y OPEX y situó el costo del primer año en torno a
S/~17~mil por intersección, dominado por la estructura mecánica y el sistema de visión. Antes de defender
una cifra definitiva deben actualizarse cotizaciones, tipo de cambio, consumo real y prorrateo de los
servicios de nube; la certificación, la obra civil y el piloto deben presupuestarse como etapas posteriores
(OE10).

\item La principal limitación del trabajo es la ausencia de fabricación, de calibración con aforos reales,
de integración con un controlador comercial certificado, de prueba vial y de auditoría de ciberseguridad.
Estas limitaciones no invalidan el diseño, pero restringen la generalización de los resultados a las
condiciones simuladas y a los supuestos documentados.

\item La arquitectura resultante constituye una base técnica trazable para un prototipo de laboratorio y un
piloto progresivo. La condición de avance debe ser demostrar cada requerimiento mandatorio en banco,
preservar la autoridad final del controlador y obtener la aprobación normativa e institucional antes de
cualquier intervención en vía pública.
\end{enumerate}

\section*{Recomendaciones}

\begin{itemize}
\item Fabricar un prototipo del módulo y verificar consumo, temperatura, respaldo, comunicaciones y
\textit{watchdog} en un banco de pruebas.
\item Probar la cámara con video real de Lima en condiciones de día, noche, contraluz, lluvia y neblina
costera, y medir la precisión de conteo, cola y velocidad.
\item Calibrar la demanda y el comportamiento vehicular de SUMO con aforos reales, ampliar el número de
semillas y publicar los archivos de configuración y de salida para garantizar la reproducibilidad.
\item Integrar el módulo a un controlador semafórico comercial en un banco \textit{hardware-in-the-loop},
sin conexión directa a las lámparas durante la etapa inicial.
\item Ejecutar pruebas autorizadas de ciberseguridad: gestión de certificados, rotación de credenciales,
restauración de respaldos y respuesta ante incidentes.
\item Revisar la normativa vial, la seguridad funcional, la compatibilidad electromagnética, la protección
ambiental y los procedimientos municipales aplicables.
\item Realizar un piloto controlado en una intersección o recinto cerrado, con plan de reversión,
supervisión humana y métricas acordadas.
\item Mantener el aprendizaje por refuerzo, la predicción CNN--LSTM y la optimización metropolitana como
líneas futuras, hasta disponer de datos, validación fuera de muestra y criterios de aceptación
(Anexo~\ref{anexo:futuros}).
\item Extender el estudio de una intersección a un corredor solo después de demostrar estabilidad local,
interoperabilidad y un beneficio estadísticamente sustentado.
\end{itemize}

% ======================================================================
%                          BIBLIOGRAFÍA
% ======================================================================
\chapter*{BIBLIOGRAFÍA}\addcontentsline{toc}{chapter}{BIBLIOGRAFÍA}\label{bibliografia}
\input{build_v3/bibliografia_items}

% ======================================================================
%                            ANEXOS
% ======================================================================
\chapter*{ANEXOS}\addcontentsline{toc}{chapter}{ANEXOS}\label{anexos}

\noindent Los archivos de programación, simulación, esquemas CAD y recursos complementarios se conservan
como material digital de soporte de la tesis. A continuación se incluyen los anexos referenciados a lo largo
del documento.

% ---------------------------------------------------------------------
\section*{ANEXO A: Reproducibilidad de la validación en SUMO}
\addcontentsline{toc}{section}{ANEXO A: Reproducibilidad de la validación en SUMO}\label{anexo:sumo}

La comparación entre el control de tiempo fijo y el control difuso adaptativo se ejecutó sobre el mismo
escenario y demanda, conservando el presupuesto de verde por ciclo, de modo que la única diferencia entre
ambas corridas es la política de reparto del tiempo de verde. Las métricas provienen de la simulación real
(no se utilizan números aleatorios sintéticos): se obtienen a nivel de vehículo mediante la interfaz de
control de SUMO. Los parámetros de reproducibilidad son los siguientes:

\begin{itemize}
\item \textbf{Motor de simulación:} SUMO mediante \texttt{libsumo}.
\item \textbf{Escenario:} Lima-Centro (\texttt{comparacion.sumocfg}); red importada de OpenStreetMap.
\item \textbf{Horizonte:} 700 pasos de simulación por corrida.
\item \textbf{Intervalo de control:} 25~s (cada intervalo se mide el estado, se calcula el ICV por
dirección y la lógica difusa recalcula el verde recomendado).
\item \textbf{Semillas:} \{1, 2, 3\}.
\item \textbf{Restricciones de seguridad:} límites duros de verde $[10, 120]$~s aplicados por una guardia de
seguridad funcional; cualquier recomendación fuera de rango se acota antes de aplicarse, registrándose el
evento.
\end{itemize}

\noindent\textbf{Tabla A.1.} Resultados por réplica (control fijo / control difuso).
\begin{center}
\begin{tabular}{@{}lcccccc@{}}
\toprule
 & \multicolumn{3}{c}{\textbf{Control fijo}} & \multicolumn{3}{c}{\textbf{Control difuso}} \\
\cmidrule(lr){2-4}\cmidrule(lr){5-7}
\textbf{Semilla} & ICV & Demora (s) & Thr. & ICV & Demora (s) & Thr. \\
\midrule
1 & 0.5039 & 496.6 & 733 & 0.4905 & 454.8 & 761 \\
2 & 0.5040 & 492.1 & 735 & 0.5049 & 515.9 & 705 \\
3 & 0.5086 & 524.3 & 707 & 0.4908 & 460.6 & 749 \\
\midrule
\textbf{Media} & 0.5055 & 504.3 & 725 & 0.4954 & 477.1 & 738 \\
\bottomrule
\end{tabular}
\end{center}
{\footnotesize\itshape Fuente: Elaboración propia a partir de la simulación en SUMO. La semilla~2 ilustra
una condición en la que el control difuso no supera al plan fijo.\par}

\vspace{0.4cm}
La discusión completa y las mejoras porcentuales agregadas se presentan en el
Capítulo~\ref{cap:validacion}.

% ---------------------------------------------------------------------
\section*{ANEXO B: Formulación de los módulos de investigación futura}
\addcontentsline{toc}{section}{ANEXO B: Módulos de investigación futura (RL y CNN--LSTM)}\label{anexo:futuros}

Los módulos de aprendizaje por refuerzo (RL) y de predicción CNN--LSTM se presentan como diseño conceptual y
\textbf{no forman parte del núcleo validado} de esta tesis: no fueron entrenados ni validados fuera de
muestra, y su desarrollo completo queda como línea futura. Se documentan aquí para sustentar la
escalabilidad de la arquitectura.

\subsection*{B.1\quad Coordinación por aprendizaje por refuerzo (A2C)}
El problema de coordinación de varias intersecciones se formula como un proceso de decisión de Markov (MDP):
\begin{itemize}
\item \textbf{Estado:} vector con el ICV por dirección de cada intersección del corredor, la longitud de
cola normalizada, la fase activa y la fracción de ciclo transcurrida.
\item \textbf{Acción:} ajuste acotado del reparto de verde y de los desfases entre intersecciones vecinas,
siempre dentro de los límites duros de seguridad.
\item \textbf{Recompensa:} combinación ponderada que penaliza la demora total y la saturación de cola y
premia el \textit{throughput} de la red.
\end{itemize}
Se propone un esquema actor--crítico de ventaja (Advantage Actor--Critic, A2C): una red actor que entrega la
política de ajuste y una red crítico que estima el valor del estado. El entrenamiento requeriría episodios
sobre escenarios SUMO calibrados, control de la distribución de datos y validación fuera de muestra antes de
cualquier uso operativo.

\subsection*{B.2\quad Predicción de congestión con CNN--LSTM}
Para anticipar la congestión se propone una arquitectura híbrida: una etapa convolucional (CNN) que extrae
patrones espaciales del estado de la red y una etapa recurrente (LSTM) que modela la evolución temporal del
ICV, entregando una predicción a un horizonte corto. Esta predicción alimentaría a la capa de coordinación
para anticipar la formación de colas. El entrenamiento y el reentrenamiento periódico exigen series
históricas representativas, que se obtendrían de la operación del propio sistema una vez desplegado.

% ---------------------------------------------------------------------
\section*{ANEXO C: Especificaciones técnicas comparativas de componentes}
\addcontentsline{toc}{section}{ANEXO C: Especificaciones técnicas de componentes}\label{anexo:componentes}

Se incluyen las tablas comparativas y las especificaciones técnicas detalladas de los componentes
seleccionados en los Capítulos~\ref{cap:diseno-electronico} y~\ref{cap:diseno-mecanico}.

\section*{C.1\quad Sensores}

\Needspace{6\baselineskip}
\captionof*{table}{\small\bfseries Tabla C-4.1. Comparativa de sensores de temperatura}

\begin{longtable}[]{@{}
  >{\raggedright\arraybackslash}p{(\columnwidth - 6\tabcolsep) * \real{0.1641}}
  >{\raggedright\arraybackslash}p{(\columnwidth - 6\tabcolsep) * \real{0.2352}}
  >{\raggedright\arraybackslash}p{(\columnwidth - 6\tabcolsep) * \real{0.3003}}
  >{\raggedright\arraybackslash}p{(\columnwidth - 6\tabcolsep) * \real{0.3003}}@{}}
\toprule
\begin{minipage}[b]{\linewidth}\raggedright
\textbf{Criterio}
\end{minipage} & \begin{minipage}[b]{\linewidth}\raggedright
\textbf{DS18B20}
\end{minipage} & \begin{minipage}[b]{\linewidth}\raggedright
\textbf{DHT22}
\end{minipage} & \begin{minipage}[b]{\linewidth}\raggedright
\textbf{LM35}
\end{minipage} \\
\begin{minipage}[b]{\linewidth}\raggedright
Tipo de salida
\end{minipage} & \begin{minipage}[b]{\linewidth}\raggedright
Digital (OneWire)
\end{minipage} & \begin{minipage}[b]{\linewidth}\raggedright
Digital (protocolo propietario)
\end{minipage} & \begin{minipage}[b]{\linewidth}\raggedright
Analógica
\end{minipage} \\
\begin{minipage}[b]{\linewidth}\raggedright
Rango de temperatura
\end{minipage} & \begin{minipage}[b]{\linewidth}\raggedright
-55°C a +125°C
\end{minipage} & \begin{minipage}[b]{\linewidth}\raggedright
-40°C a +80°C
\end{minipage} & \begin{minipage}[b]{\linewidth}\raggedright
-55°C a +150°C
\end{minipage} \\
\begin{minipage}[b]{\linewidth}\raggedright
Precisión
\end{minipage} & \begin{minipage}[b]{\linewidth}\raggedright
\ensuremath{\pm}0.5°C (-10°C a +85°C)
\end{minipage} & \begin{minipage}[b]{\linewidth}\raggedright
\ensuremath{\pm}0.5°C (15°C a 25°C), \ensuremath{\pm}2°C (fuera del rango)
\end{minipage} & \begin{minipage}[b]{\linewidth}\raggedright
\ensuremath{\pm}0.5°C (a 25°C)
\end{minipage} \\
\midrule
\endhead
\bottomrule
\endlastfoot
\end{longtable}

\Needspace{6\baselineskip}
\captionof*{table}{\small\bfseries Tabla C-4.2. Especificaciones técnicas del Sensor DS18B20}

\begin{longtable}[]{@{}
  >{\raggedright\arraybackslash}p{(\columnwidth - 2\tabcolsep) * \real{0.3005}}
  >{\raggedright\arraybackslash}p{(\columnwidth - 2\tabcolsep) * \real{0.6995}}@{}}
\toprule
\begin{minipage}[b]{\linewidth}\raggedright
\textbf{Característica}
\end{minipage} & \begin{minipage}[b]{\linewidth}\raggedright
\textbf{Descripción}
\end{minipage} \\
\begin{minipage}[b]{\linewidth}\raggedright
\textbf{Función}
\end{minipage} & \begin{minipage}[b]{\linewidth}\raggedright
Medición digital de temperatura de 9-bit a 12-bit en Celsius
\end{minipage} \\
\begin{minipage}[b]{\linewidth}\raggedright
\textbf{Protocolo comunicación}
\end{minipage} & \begin{minipage}[b]{\linewidth}\raggedright
1-Wire bus
\end{minipage} \\
\begin{minipage}[b]{\linewidth}\raggedright
\textbf{Rango de temperatura}
\end{minipage} & \begin{minipage}[b]{\linewidth}\raggedright
-55°C to +125°C (-67°F to +257°F)
\end{minipage} \\
\begin{minipage}[b]{\linewidth}\raggedright
\textbf{Precisión}
\end{minipage} & \begin{minipage}[b]{\linewidth}\raggedright
\ensuremath{\pm}5°C (de -10°C a +85°C)
\end{minipage} \\
\midrule
\endhead
\bottomrule
\endlastfoot
\end{longtable}

\Needspace{6\baselineskip}
\captionof*{table}{\small\bfseries Tabla C-4.3. Comparativa de sensores de corriente}

\begin{longtable}[]{@{}
  >{\raggedright\arraybackslash}p{(\columnwidth - 6\tabcolsep) * \real{0.2190}}
  >{\raggedright\arraybackslash}p{(\columnwidth - 6\tabcolsep) * \real{0.1796}}
  >{\raggedright\arraybackslash}p{(\columnwidth - 6\tabcolsep) * \real{0.3007}}
  >{\raggedright\arraybackslash}p{(\columnwidth - 6\tabcolsep) * \real{0.3007}}@{}}
\toprule
\begin{minipage}[b]{\linewidth}\raggedright
\textbf{Criterio}
\end{minipage} & \begin{minipage}[b]{\linewidth}\raggedright
\textbf{WCS1600 (Efecto Hall)}
\end{minipage} & \begin{minipage}[b]{\linewidth}\raggedright
\textbf{Sensor Shunt (ACS712)}
\end{minipage} & \begin{minipage}[b]{\linewidth}\raggedright
\textbf{Transformador de Corriente (CT)}
\end{minipage} \\
\begin{minipage}[b]{\linewidth}\raggedright
Principio de funcionamiento
\end{minipage} & \begin{minipage}[b]{\linewidth}\raggedright
Efecto Hall
\end{minipage} & \begin{minipage}[b]{\linewidth}\raggedright
Resistencia shunt + Hall
\end{minipage} & \begin{minipage}[b]{\linewidth}\raggedright
Inducción magnética
\end{minipage} \\
\begin{minipage}[b]{\linewidth}\raggedright
Rango de medición
\end{minipage} & \begin{minipage}[b]{\linewidth}\raggedright
\ensuremath{\pm}100A DC, 70A RMS AC
\end{minipage} & \begin{minipage}[b]{\linewidth}\raggedright
\ensuremath{\pm}5A, \ensuremath{\pm}20A, \ensuremath{\pm}30A
\end{minipage} & \begin{minipage}[b]{\linewidth}\raggedright
0-100A AC
\end{minipage} \\
\begin{minipage}[b]{\linewidth}\raggedright
Salida
\end{minipage} & \begin{minipage}[b]{\linewidth}\raggedright
Analógica (0.3-4.7V)
\end{minipage} & \begin{minipage}[b]{\linewidth}\raggedright
Analógica (0-5V)
\end{minipage} & \begin{minipage}[b]{\linewidth}\raggedright
Analógica (AC)
\end{minipage} \\
\midrule
\endhead
\bottomrule
\endlastfoot
\end{longtable}

\Needspace{6\baselineskip}
\captionof*{table}{\small\bfseries Tabla C-4.4. Especificaciones técnicas del Sensor WCS1600}

\begin{longtable}[]{@{}
  >{\raggedright\arraybackslash}p{(\columnwidth - 2\tabcolsep) * \real{0.3845}}
  >{\raggedright\arraybackslash}p{(\columnwidth - 2\tabcolsep) * \real{0.6155}}@{}}
\toprule
\begin{minipage}[b]{\linewidth}\raggedright
\textbf{Característica}
\end{minipage} & \begin{minipage}[b]{\linewidth}\raggedright
\textbf{Descripción}
\end{minipage} \\
\begin{minipage}[b]{\linewidth}\raggedright
\textbf{Función}
\end{minipage} & \begin{minipage}[b]{\linewidth}\raggedright
Medición de corriente DC o AC utilizando el efecto Hall
\end{minipage} \\
\begin{minipage}[b]{\linewidth}\raggedright
\textbf{Rango de medición}
\end{minipage} & \begin{minipage}[b]{\linewidth}\raggedright
\ensuremath{\pm}100A (DC), 70A RMS(AC)
\end{minipage} \\
\begin{minipage}[b]{\linewidth}\raggedright
\textbf{Temperatura de funcionamiento}
\end{minipage} & \begin{minipage}[b]{\linewidth}\raggedright
-20°C \textasciitilde{} 125°C
\end{minipage} \\
\begin{minipage}[b]{\linewidth}\raggedright
\textbf{Diámetro máximo del conductor}
\end{minipage} & \begin{minipage}[b]{\linewidth}\raggedright
9mm
\end{minipage} \\
\midrule
\endhead
\bottomrule
\endlastfoot
\end{longtable}

\Needspace{6\baselineskip}
\captionof*{table}{\small\bfseries Tabla C-4.5. Comparativa de sensores de presencia}

\begin{longtable}[]{@{}
  >{\raggedright\arraybackslash}p{(\columnwidth - 6\tabcolsep) * \real{0.1908}}
  >{\raggedright\arraybackslash}p{(\columnwidth - 6\tabcolsep) * \real{0.1899}}
  >{\raggedright\arraybackslash}p{(\columnwidth - 6\tabcolsep) * \real{0.2628}}
  >{\raggedright\arraybackslash}p{(\columnwidth - 6\tabcolsep) * \real{0.3565}}@{}}
\toprule
\begin{minipage}[b]{\linewidth}\raggedright
\textbf{Criterio}
\end{minipage} & \begin{minipage}[b]{\linewidth}\raggedright
\textbf{URM06 (Ultrasónico)}
\end{minipage} & \begin{minipage}[b]{\linewidth}\raggedright
\textbf{HC-SR04 (Ultrasónico)}
\end{minipage} & \begin{minipage}[b]{\linewidth}\raggedright
\textbf{Sensor Infrarrojo (PIR/IR)}
\end{minipage} \\
\begin{minipage}[b]{\linewidth}\raggedright
Principio de funcionamiento
\end{minipage} & \begin{minipage}[b]{\linewidth}\raggedright
Ultrasonido 40kHz
\end{minipage} & \begin{minipage}[b]{\linewidth}\raggedright
Ultrasonido 40kHz
\end{minipage} & \begin{minipage}[b]{\linewidth}\raggedright
Detección infrarroja
\end{minipage} \\
\begin{minipage}[b]{\linewidth}\raggedright
Rango de detección
\end{minipage} & \begin{minipage}[b]{\linewidth}\raggedright
20cm a 10m
\end{minipage} & \begin{minipage}[b]{\linewidth}\raggedright
2cm a 4m
\end{minipage} & \begin{minipage}[b]{\linewidth}\raggedright
3m a 7m (típico)
\end{minipage} \\
\begin{minipage}[b]{\linewidth}\raggedright
Ángulo del haz
\end{minipage} & \begin{minipage}[b]{\linewidth}\raggedright
15° (haz estrecho)
\end{minipage} & \begin{minipage}[b]{\linewidth}\raggedright
30° (haz amplio)
\end{minipage} & \begin{minipage}[b]{\linewidth}\raggedright
110° (amplio)
\end{minipage} \\
\midrule
\endhead
\bottomrule
\endlastfoot
\end{longtable}

\Needspace{6\baselineskip}
\captionof*{table}{\small\bfseries Tabla C-4.6. Especificaciones técnicas del Sensor URM06}

\begin{longtable}[]{@{}
  >{\raggedright\arraybackslash}p{(\columnwidth - 2\tabcolsep) * \real{0.3073}}
  >{\raggedright\arraybackslash}p{(\columnwidth - 2\tabcolsep) * \real{0.6927}}@{}}
\toprule
\begin{minipage}[b]{\linewidth}\raggedright
\textbf{Característica}
\end{minipage} & \begin{minipage}[b]{\linewidth}\raggedright
\textbf{Descripción}
\end{minipage} \\
\begin{minipage}[b]{\linewidth}\raggedright
\textbf{Función}
\end{minipage} & \begin{minipage}[b]{\linewidth}\raggedright
Detección ultrasónica de presencia/distancia de corto a largo alcance
\end{minipage} \\
\begin{minipage}[b]{\linewidth}\raggedright
\textbf{Tipo}
\end{minipage} & \begin{minipage}[b]{\linewidth}\raggedright
Sensor ultrasónico de haz estrecho
\end{minipage} \\
\begin{minipage}[b]{\linewidth}\raggedright
\textbf{Rango de detección}
\end{minipage} & \begin{minipage}[b]{\linewidth}\raggedright
20cm \textasciitilde{} 10m (1000cm)
\end{minipage} \\
\midrule
\endhead
\bottomrule
\endlastfoot
\end{longtable}

\Needspace{6\baselineskip}
\captionof*{table}{\small\bfseries Tabla C-4.7. Comparativa de sensores de presencia}

\begin{longtable}[]{@{}
  >{\raggedright\arraybackslash}p{(\columnwidth - 8\tabcolsep) * \real{0.1991}}
  >{\raggedright\arraybackslash}p{(\columnwidth - 8\tabcolsep) * \real{0.2301}}
  >{\raggedright\arraybackslash}p{(\columnwidth - 8\tabcolsep) * \real{0.2177}}
  >{\raggedright\arraybackslash}p{(\columnwidth - 8\tabcolsep) * \real{0.1696}}
  >{\raggedright\arraybackslash}p{(\columnwidth - 8\tabcolsep) * \real{0.1836}}@{}}
\toprule
\begin{minipage}[b]{\linewidth}\raggedright
\textbf{Criterio}
\end{minipage} & \begin{minipage}[b]{\linewidth}\raggedright
\textbf{DS-2CE57H0T-VPITF (Domo)}
\end{minipage} & \begin{minipage}[b]{\linewidth}\raggedright
\textbf{Cámara PTZ Motorizada}
\end{minipage} & \begin{minipage}[b]{\linewidth}\raggedright
\textbf{Cámara IP Estándar}
\end{minipage} & \begin{minipage}[b]{\linewidth}\raggedright
\textbf{Cámara Bullet}
\end{minipage} \\
\begin{minipage}[b]{\linewidth}\raggedright
Tipo de montaje
\end{minipage} & \begin{minipage}[b]{\linewidth}\raggedright
Domo fijo
\end{minipage} & \begin{minipage}[b]{\linewidth}\raggedright
PTZ (Pan-Tilt-Zoom)
\end{minipage} & \begin{minipage}[b]{\linewidth}\raggedright
Fija
\end{minipage} & \begin{minipage}[b]{\linewidth}\raggedright
Fija cilíndrica
\end{minipage} \\
\begin{minipage}[b]{\linewidth}\raggedright
Frame rate
\end{minipage} & \begin{minipage}[b]{\linewidth}\raggedright
20 fps
\end{minipage} & \begin{minipage}[b]{\linewidth}\raggedright
25-30 fps
\end{minipage} & \begin{minipage}[b]{\linewidth}\raggedright
25-30 fps
\end{minipage} & \begin{minipage}[b]{\linewidth}\raggedright
25-30 fps
\end{minipage} \\
\begin{minipage}[b]{\linewidth}\raggedright
Visión nocturna
\end{minipage} & \begin{minipage}[b]{\linewidth}\raggedright
EXIR 2.0 (20m)
\end{minipage} & \begin{minipage}[b]{\linewidth}\raggedright
IR estándar (30m)
\end{minipage} & \begin{minipage}[b]{\linewidth}\raggedright
Varía
\end{minipage} & \begin{minipage}[b]{\linewidth}\raggedright
IR estándar (20-30m)
\end{minipage} \\
\begin{minipage}[b]{\linewidth}\raggedright
Protección ambiental
\end{minipage} & \begin{minipage}[b]{\linewidth}\raggedright
IP67 + IK10
\end{minipage} & \begin{minipage}[b]{\linewidth}\raggedright
IP66 + IK10
\end{minipage} & \begin{minipage}[b]{\linewidth}\raggedright
IP65 (típica)
\end{minipage} & \begin{minipage}[b]{\linewidth}\raggedright
IP66
\end{minipage} \\
\begin{minipage}[b]{\linewidth}\raggedright
Resistencia vandalismo
\end{minipage} & \begin{minipage}[b]{\linewidth}\raggedright
Excelente (IK10)
\end{minipage} & \begin{minipage}[b]{\linewidth}\raggedright
Buena
\end{minipage} & \begin{minipage}[b]{\linewidth}\raggedright
Limitada
\end{minipage} & \begin{minipage}[b]{\linewidth}\raggedright
Moderada
\end{minipage} \\
\begin{minipage}[b]{\linewidth}\raggedright
Consumo energía
\end{minipage} & \begin{minipage}[b]{\linewidth}\raggedright
4W
\end{minipage} & \begin{minipage}[b]{\linewidth}\raggedright
15-30W
\end{minipage} & \begin{minipage}[b]{\linewidth}\raggedright
5-8W
\end{minipage} & \begin{minipage}[b]{\linewidth}\raggedright
4-6W
\end{minipage} \\
\midrule
\endhead
\bottomrule
\endlastfoot
\end{longtable}

\Needspace{6\baselineskip}
\captionof*{table}{\small\bfseries Tabla C-4.8. Especificaciones técnicas del DS-2CE57H0T-VPITF}

\begin{longtable}[]{@{}
  >{\raggedright\arraybackslash}p{(\columnwidth - 2\tabcolsep) * \real{0.3065}}
  >{\raggedright\arraybackslash}p{(\columnwidth - 2\tabcolsep) * \real{0.6935}}@{}}
\toprule
\begin{minipage}[b]{\linewidth}\raggedright
\textbf{Característica}
\end{minipage} & \begin{minipage}[b]{\linewidth}\raggedright
\textbf{Descripción}
\end{minipage} \\
\begin{minipage}[b]{\linewidth}\raggedright
\textbf{Función}
\end{minipage} & \begin{minipage}[b]{\linewidth}\raggedright
Videovigilancia domo fijo de 5 MP con visión nocturna
\end{minipage} \\
\begin{minipage}[b]{\linewidth}\raggedright
\textbf{Resolución máxima}
\end{minipage} & \begin{minipage}[b]{\linewidth}\raggedright
2560 (H) \ensuremath{\times} 1944 (V)
\end{minipage} \\
\begin{minipage}[b]{\linewidth}\raggedright
\textbf{Frame rate}
\end{minipage} & \begin{minipage}[b]{\linewidth}\raggedright
20 fps
\end{minipage} \\
\begin{minipage}[b]{\linewidth}\raggedright
\textbf{Protección}
\end{minipage} & \begin{minipage}[b]{\linewidth}\raggedright
Resistente al agua y polvo (IP67) y prueba de vandalismo (IK10)
\end{minipage} \\
\midrule
\endhead
\bottomrule
\endlastfoot
\end{longtable}

\section*{C.2\quad Actuadores y drivers}

\Needspace{6\baselineskip}
\captionof*{table}{\small\bfseries Tabla C-4.9. Comparativa de actuadores para movilización de cámara}

\begin{longtable}[]{@{}
  >{\raggedright\arraybackslash}p{(\columnwidth - 6\tabcolsep) * \real{0.2102}}
  >{\raggedright\arraybackslash}p{(\columnwidth - 6\tabcolsep) * \real{0.2603}}
  >{\raggedright\arraybackslash}p{(\columnwidth - 6\tabcolsep) * \real{0.3321}}
  >{\raggedright\arraybackslash}p{(\columnwidth - 6\tabcolsep) * \real{0.1973}}@{}}
\toprule
\begin{minipage}[b]{\linewidth}\raggedright
\textbf{Criterio}
\end{minipage} & \begin{minipage}[b]{\linewidth}\raggedright
\textbf{Motor Paso a Paso}
\end{minipage} & \begin{minipage}[b]{\linewidth}\raggedright
\textbf{Motor DC con Encoders}
\end{minipage} & \begin{minipage}[b]{\linewidth}\raggedright
\textbf{Servomotor}
\end{minipage} \\
\begin{minipage}[b]{\linewidth}\raggedright
Control de posición
\end{minipage} & \begin{minipage}[b]{\linewidth}\raggedright
Lazo abierto (sin retroalimentación)
\end{minipage} & \begin{minipage}[b]{\linewidth}\raggedright
Lazo cerrado (requiere encoders)
\end{minipage} & \begin{minipage}[b]{\linewidth}\raggedright
Lazo cerrado (interno)
\end{minipage} \\
\begin{minipage}[b]{\linewidth}\raggedright
Precisión angular
\end{minipage} & \begin{minipage}[b]{\linewidth}\raggedright
1.8° por paso (200 pasos/rev)
\end{minipage} & \begin{minipage}[b]{\linewidth}\raggedright
Depende del encoder
\end{minipage} & \begin{minipage}[b]{\linewidth}\raggedright
0.1-0.5° (típico)
\end{minipage} \\
\begin{minipage}[b]{\linewidth}\raggedright
Mantenimiento de posición
\end{minipage} & \begin{minipage}[b]{\linewidth}\raggedright
Excelente sin energía adicional
\end{minipage} & \begin{minipage}[b]{\linewidth}\raggedright
Requiere consumo continuo
\end{minipage} & \begin{minipage}[b]{\linewidth}\raggedright
Requiere señal PWM continua
\end{minipage} \\
\midrule
\endhead
\bottomrule
\endlastfoot
\end{longtable}

\Needspace{6\baselineskip}
\captionof*{table}{\small\bfseries Tabla C-4.10. Especificaciones técnicas del Motor PaP Stepper Nema 23}

\begin{longtable}[]{@{}
  >{\raggedright\arraybackslash}p{(\columnwidth - 2\tabcolsep) * \real{0.3162}}
  >{\raggedright\arraybackslash}p{(\columnwidth - 2\tabcolsep) * \real{0.6838}}@{}}
\toprule
\begin{minipage}[b]{\linewidth}\raggedright
\textbf{Característica}
\end{minipage} & \begin{minipage}[b]{\linewidth}\raggedright
\textbf{Descripción}
\end{minipage} \\
\begin{minipage}[b]{\linewidth}\raggedright
\textbf{Corriente nominal por fase}
\end{minipage} & \begin{minipage}[b]{\linewidth}\raggedright
4.2A
\end{minipage} \\
\begin{minipage}[b]{\linewidth}\raggedright
\textbf{Torque máximo detenido}
\end{minipage} & \begin{minipage}[b]{\linewidth}\raggedright
18.4 kg.cm (1.8 N.m)
\end{minipage} \\
\begin{minipage}[b]{\linewidth}\raggedright
\textbf{Diámetro del eje}
\end{minipage} & \begin{minipage}[b]{\linewidth}\raggedright
8 mm (Eje tipo D)
\end{minipage} \\
\begin{minipage}[b]{\linewidth}\raggedright
\textbf{Largo del eje}
\end{minipage} & \begin{minipage}[b]{\linewidth}\raggedright
21 mm
\end{minipage} \\
\begin{minipage}[b]{\linewidth}\raggedright
\textbf{Perfil NEMA 23}
\end{minipage} & \begin{minipage}[b]{\linewidth}\raggedright
56*56 mm
\end{minipage} \\
\begin{minipage}[b]{\linewidth}\raggedright
\textbf{Agujeros para montaje}
\end{minipage} & \begin{minipage}[b]{\linewidth}\raggedright
4*D5.2 mm pasante
\end{minipage} \\
\midrule
\endhead
\bottomrule
\endlastfoot
\end{longtable}

\Needspace{6\baselineskip}
\captionof*{table}{\small\bfseries Tabla C-4.11. Driver DM542T - Especificaciones Técnicas}

\begin{longtable}[]{@{}
  >{\raggedright\arraybackslash}p{(\columnwidth - 2\tabcolsep) * \real{0.3309}}
  >{\raggedright\arraybackslash}p{(\columnwidth - 2\tabcolsep) * \real{0.6691}}@{}}
\toprule
\begin{minipage}[b]{\linewidth}\raggedright
\textbf{Característica}
\end{minipage} & \begin{minipage}[b]{\linewidth}\raggedright
\textbf{Descripción}
\end{minipage} \\
\begin{minipage}[b]{\linewidth}\raggedright
\textbf{Voltaje de entrada}
\end{minipage} & \begin{minipage}[b]{\linewidth}\raggedright
18V - 50V DC
\end{minipage} \\
\begin{minipage}[b]{\linewidth}\raggedright
\textbf{Corriente de salida}
\end{minipage} & \begin{minipage}[b]{\linewidth}\raggedright
1.0A - 4.5A
\end{minipage} \\
\begin{minipage}[b]{\linewidth}\raggedright
\textbf{Corriente máxima RMS}
\end{minipage} & \begin{minipage}[b]{\linewidth}\raggedright
3.2A
\end{minipage} \\
\begin{minipage}[b]{\linewidth}\raggedright
\textbf{Resolución de micropasos}
\end{minipage} & \begin{minipage}[b]{\linewidth}\raggedright
15 opciones: desde 400 hasta 25,600 pasos/revolución
\end{minipage} \\
\begin{minipage}[b]{\linewidth}\raggedright
\textbf{Temperatura de operación}
\end{minipage} & \begin{minipage}[b]{\linewidth}\raggedright
-10°C a +45°C
\end{minipage} \\
\midrule
\endhead
\bottomrule
\endlastfoot
\end{longtable}

\Needspace{6\baselineskip}
\captionof*{table}{\small\bfseries Tabla C-4.12. Conexiones Físicas y Configuración}

\begin{longtable}[]{@{}
  >{\raggedright\arraybackslash}p{(\columnwidth - 2\tabcolsep) * \real{0.3251}}
  >{\raggedright\arraybackslash}p{(\columnwidth - 2\tabcolsep) * \real{0.6749}}@{}}
\toprule
\begin{minipage}[b]{\linewidth}\raggedright
\textbf{Parámetro}
\end{minipage} & \begin{minipage}[b]{\linewidth}\raggedright
\textbf{Especificación}
\end{minipage} \\
\begin{minipage}[b]{\linewidth}\raggedright
\textbf{Conexión física}
\end{minipage} & \begin{minipage}[b]{\linewidth}\raggedright
GPIO 40 pines + Header apilable incluido
\end{minipage} \\
\begin{minipage}[b]{\linewidth}\raggedright
\textbf{Alimentación}
\end{minipage} & \begin{minipage}[b]{\linewidth}\raggedright
5V vía GPIO o USB-C externo (hasta 27W)
\end{minipage} \\
\begin{minipage}[b]{\linewidth}\raggedright
\textbf{Comunicación con RPi 5}
\end{minipage} & \begin{minipage}[b]{\linewidth}\raggedright
USB 2.0
\end{minipage} \\
\begin{minipage}[b]{\linewidth}\raggedright
\textbf{Control GPIO}
\end{minipage} & \begin{minipage}[b]{\linewidth}\raggedright
RST (GPIO 29), PWR (GPIO 31), RX/TX configurable
\end{minipage} \\
\begin{minipage}[b]{\linewidth}\raggedright
\textbf{Ranura SIM}
\end{minipage} & \begin{minipage}[b]{\linewidth}\raggedright
Nano SIM 3V/1.8V
\end{minipage} \\
\begin{minipage}[b]{\linewidth}\raggedright
\textbf{Antenas}
\end{minipage} & \begin{minipage}[b]{\linewidth}\raggedright
2x LTE + 1x GPS (conectores IPEX-1)
\end{minipage} \\
\midrule
\endhead
\bottomrule
\endlastfoot
\end{longtable}

\section*{C.3\quad Alimentación y protecciones}

\Needspace{6\baselineskip}
\captionof*{table}{\small\bfseries Tabla C-4.16. Especificaciones técnicas SMPS}

\begin{longtable}[]{@{}
  >{\raggedright\arraybackslash}p{(\columnwidth - 2\tabcolsep) * \real{0.3919}}
  >{\raggedright\arraybackslash}p{(\columnwidth - 2\tabcolsep) * \real{0.6081}}@{}}
\toprule
\begin{minipage}[b]{\linewidth}\raggedright
\textbf{Característica}
\end{minipage} & \begin{minipage}[b]{\linewidth}\raggedright
\textbf{Especificación}
\end{minipage} \\
\begin{minipage}[b]{\linewidth}\raggedright
\textbf{Salida de Voltaje}
\end{minipage} & \begin{minipage}[b]{\linewidth}\raggedright
12 VDC
\end{minipage} \\
\begin{minipage}[b]{\linewidth}\raggedright
\textbf{Corriente de salida}
\end{minipage} & \begin{minipage}[b]{\linewidth}\raggedright
8 A
\end{minipage} \\
\begin{minipage}[b]{\linewidth}\raggedright
\textbf{Potencia máxima}
\end{minipage} & \begin{minipage}[b]{\linewidth}\raggedright
96 W
\end{minipage} \\
\begin{minipage}[b]{\linewidth}\raggedright
\textbf{Entrada de Voltaje}
\end{minipage} & \begin{minipage}[b]{\linewidth}\raggedright
24 VDC
\end{minipage} \\
\begin{minipage}[b]{\linewidth}\raggedright
\textbf{Eficiencia típica}
\end{minipage} & \begin{minipage}[b]{\linewidth}\raggedright
88.2 \%
\end{minipage} \\
\begin{minipage}[b]{\linewidth}\raggedright
\textbf{Protecciones}
\end{minipage} & \begin{minipage}[b]{\linewidth}\raggedright
Cortocircuitos, sobrecarga, sobretensión
\end{minipage} \\
\begin{minipage}[b]{\linewidth}\raggedright
\textbf{Dimensiones}
\end{minipage} & \begin{minipage}[b]{\linewidth}\raggedright
32 x 124 x 102 mm
\end{minipage} \\
\begin{minipage}[b]{\linewidth}\raggedright
\textbf{Peso}
\end{minipage} & \begin{minipage}[b]{\linewidth}\raggedright
425 g
\end{minipage} \\
\midrule
\endhead
\bottomrule
\endlastfoot
\end{longtable}

\Needspace{6\baselineskip}
\captionof*{table}{\small\bfseries Tabla C-4.17. Especificaciones técnicas SMPS}

\begin{longtable}[]{@{}
  >{\raggedright\arraybackslash}p{(\columnwidth - 2\tabcolsep) * \real{0.3919}}
  >{\raggedright\arraybackslash}p{(\columnwidth - 2\tabcolsep) * \real{0.6081}}@{}}
\toprule
\begin{minipage}[b]{\linewidth}\raggedright
\textbf{Característica}
\end{minipage} & \begin{minipage}[b]{\linewidth}\raggedright
\textbf{Especificación}
\end{minipage} \\
\begin{minipage}[b]{\linewidth}\raggedright
\textbf{Salida de Voltaje}
\end{minipage} & \begin{minipage}[b]{\linewidth}\raggedright
5 VDC
\end{minipage} \\
\begin{minipage}[b]{\linewidth}\raggedright
\textbf{Corriente de salida}
\end{minipage} & \begin{minipage}[b]{\linewidth}\raggedright
10 A
\end{minipage} \\
\begin{minipage}[b]{\linewidth}\raggedright
\textbf{Potencia máxima}
\end{minipage} & \begin{minipage}[b]{\linewidth}\raggedright
50 W
\end{minipage} \\
\begin{minipage}[b]{\linewidth}\raggedright
\textbf{Entrada de Voltaje}
\end{minipage} & \begin{minipage}[b]{\linewidth}\raggedright
24 VDC
\end{minipage} \\
\begin{minipage}[b]{\linewidth}\raggedright
\textbf{Eficiencia típica}
\end{minipage} & \begin{minipage}[b]{\linewidth}\raggedright
81,5\%
\end{minipage} \\
\begin{minipage}[b]{\linewidth}\raggedright
\textbf{Protecciones}
\end{minipage} & \begin{minipage}[b]{\linewidth}\raggedright
Cortocircuitos, sobrecarga, sobretensión
\end{minipage} \\
\begin{minipage}[b]{\linewidth}\raggedright
\textbf{Peso}
\end{minipage} & \begin{minipage}[b]{\linewidth}\raggedright
425 g
\end{minipage} \\
\midrule
\endhead
\bottomrule
\endlastfoot
\end{longtable}

\Needspace{6\baselineskip}
\captionof*{table}{\small\bfseries Tabla C-4.18. Especificaciones técnicas SMPS}

\begin{longtable}[]{@{}
  >{\raggedright\arraybackslash}p{(\columnwidth - 2\tabcolsep) * \real{0.3919}}
  >{\raggedright\arraybackslash}p{(\columnwidth - 2\tabcolsep) * \real{0.6081}}@{}}
\toprule
\begin{minipage}[b]{\linewidth}\raggedright
\textbf{Parámetro}
\end{minipage} & \begin{minipage}[b]{\linewidth}\raggedright
\textbf{Especificación}
\end{minipage} \\
\begin{minipage}[b]{\linewidth}\raggedright
\textbf{Modelo}
\end{minipage} & \begin{minipage}[b]{\linewidth}\raggedright
S-360-24
\end{minipage} \\
\begin{minipage}[b]{\linewidth}\raggedright
\textbf{Potencia nominal}
\end{minipage} & \begin{minipage}[b]{\linewidth}\raggedright
360W
\end{minipage} \\
\begin{minipage}[b]{\linewidth}\raggedright
\textbf{Voltaje entrada}
\end{minipage} & \begin{minipage}[b]{\linewidth}\raggedright
220V AC, 60Hz
\end{minipage} \\
\begin{minipage}[b]{\linewidth}\raggedright
\textbf{Voltaje salida principal}
\end{minipage} & \begin{minipage}[b]{\linewidth}\raggedright
24V DC
\end{minipage} \\
\begin{minipage}[b]{\linewidth}\raggedright
\textbf{Corriente máxima}
\end{minipage} & \begin{minipage}[b]{\linewidth}\raggedright
15A
\end{minipage} \\
\begin{minipage}[b]{\linewidth}\raggedright
\textbf{Eficiencia}
\end{minipage} & \begin{minipage}[b]{\linewidth}\raggedright
\textgreater85\%
\end{minipage} \\
\begin{minipage}[b]{\linewidth}\raggedright
\textbf{Regulación de línea}
\end{minipage} & \begin{minipage}[b]{\linewidth}\raggedright
\ensuremath{\pm}0.5\%
\end{minipage} \\
\begin{minipage}[b]{\linewidth}\raggedright
\textbf{Regulación de carga}
\end{minipage} & \begin{minipage}[b]{\linewidth}\raggedright
\ensuremath{\pm}1\%
\end{minipage} \\
\begin{minipage}[b]{\linewidth}\raggedright
\textbf{Temperatura operativa}
\end{minipage} & \begin{minipage}[b]{\linewidth}\raggedright
-10°C a +50°C
\end{minipage} \\
\begin{minipage}[b]{\linewidth}\raggedright
\textbf{Dimensiones}
\end{minipage} & \begin{minipage}[b]{\linewidth}\raggedright
215 \ensuremath{\times} 115 \ensuremath{\times} 51 mm
\end{minipage} \\
\midrule
\endhead
\bottomrule
\endlastfoot
\end{longtable}

\Needspace{6\baselineskip}
\captionof*{table}{\small\bfseries Tabla C-4.19. Especificaciones del UPS}

\begin{longtable}[]{@{}
  >{\raggedright\arraybackslash}p{(\columnwidth - 2\tabcolsep) * \real{0.3796}}
  >{\raggedright\arraybackslash}p{(\columnwidth - 2\tabcolsep) * \real{0.6204}}@{}}
\toprule
\begin{minipage}[b]{\linewidth}\raggedright
\textbf{Parámetro}
\end{minipage} & \begin{minipage}[b]{\linewidth}\raggedright
\textbf{Especificación}
\end{minipage} \\
\begin{minipage}[b]{\linewidth}\raggedright
\textbf{Potencia nominal}
\end{minipage} & \begin{minipage}[b]{\linewidth}\raggedright
1500VA / 1000W
\end{minipage} \\
\begin{minipage}[b]{\linewidth}\raggedright
\textbf{Voltaje entrada}
\end{minipage} & \begin{minipage}[b]{\linewidth}\raggedright
160-286V AC
\end{minipage} \\
\begin{minipage}[b]{\linewidth}\raggedright
\textbf{Voltaje salida}
\end{minipage} & \begin{minipage}[b]{\linewidth}\raggedright
230V AC \ensuremath{\pm} 3\%
\end{minipage} \\
\begin{minipage}[b]{\linewidth}\raggedright
\textbf{Frecuencia}
\end{minipage} & \begin{minipage}[b]{\linewidth}\raggedright
50/60 Hz \ensuremath{\pm} 3Hz
\end{minipage} \\
\begin{minipage}[b]{\linewidth}\raggedright
\textbf{Forma de onda}
\end{minipage} & \begin{minipage}[b]{\linewidth}\raggedright
Senoidal pura
\end{minipage} \\
\begin{minipage}[b]{\linewidth}\raggedright
\textbf{Baterías}
\end{minipage} & \begin{minipage}[b]{\linewidth}\raggedright
2 \ensuremath{\times} 12V 12Ah AGM
\end{minipage} \\
\begin{minipage}[b]{\linewidth}\raggedright
\textbf{Protecciones}
\end{minipage} & \begin{minipage}[b]{\linewidth}\raggedright
Sobrecarga, cortocircuito, sobretensión
\end{minipage} \\
\midrule
\endhead
\bottomrule
\endlastfoot
\end{longtable}

\Needspace{6\baselineskip}
\captionof*{table}{\small\bfseries Tabla C-4.20. Especificaciones del Breaker}

\begin{longtable}[]{@{}
  >{\raggedright\arraybackslash}p{(\columnwidth - 2\tabcolsep) * \real{0.5335}}
  >{\raggedright\arraybackslash}p{(\columnwidth - 2\tabcolsep) * \real{0.4665}}@{}}
\toprule
\begin{minipage}[b]{\linewidth}\raggedright
\textbf{Parámetro}
\end{minipage} & \begin{minipage}[b]{\linewidth}\raggedright
\textbf{Especificación}
\end{minipage} \\
\begin{minipage}[b]{\linewidth}\raggedright
\textbf{Corriente nominal}
\end{minipage} & \begin{minipage}[b]{\linewidth}\raggedright
6A
\end{minipage} \\
\begin{minipage}[b]{\linewidth}\raggedright
\textbf{Curva de disparo}
\end{minipage} & \begin{minipage}[b]{\linewidth}\raggedright
Tipo C (5-10\ensuremath{\times}In)
\end{minipage} \\
\begin{minipage}[b]{\linewidth}\raggedright
\textbf{Capacidad de ruptura}
\end{minipage} & \begin{minipage}[b]{\linewidth}\raggedright
6kA a 230/400V AC
\end{minipage} \\
\begin{minipage}[b]{\linewidth}\raggedright
\textbf{Polos}
\end{minipage} & \begin{minipage}[b]{\linewidth}\raggedright
1 polo
\end{minipage} \\
\midrule
\endhead
\bottomrule
\endlastfoot
\end{longtable}

\Needspace{6\baselineskip}
\captionof*{table}{\small\bfseries Tabla C-4.21. Especificaciones del DPS}

\begin{longtable}[]{@{}
  >{\raggedright\arraybackslash}p{(\columnwidth - 2\tabcolsep) * \real{0.5948}}
  >{\raggedright\arraybackslash}p{(\columnwidth - 2\tabcolsep) * \real{0.4052}}@{}}
\toprule
\begin{minipage}[b]{\linewidth}\raggedright
\textbf{Parámetro}
\end{minipage} & \begin{minipage}[b]{\linewidth}\raggedright
\textbf{Especificación}
\end{minipage} \\
\begin{minipage}[b]{\linewidth}\raggedright
\textbf{Voltaje nominal}
\end{minipage} & \begin{minipage}[b]{\linewidth}\raggedright
\textbf{230V AC}
\end{minipage} \\
\begin{minipage}[b]{\linewidth}\raggedright
\textbf{Corriente de descarga}
\end{minipage} & \begin{minipage}[b]{\linewidth}\raggedright
\textbf{10kA (8/20\ensuremath{\mu}s)}
\end{minipage} \\
\begin{minipage}[b]{\linewidth}\raggedright
\textbf{Nivel de protección}
\end{minipage} & \begin{minipage}[b]{\linewidth}\raggedright
\textbf{\ensuremath{\leq} 1.5kV}
\end{minipage} \\
\begin{minipage}[b]{\linewidth}\raggedright
\textbf{Tiempo de respuesta}
\end{minipage} & \begin{minipage}[b]{\linewidth}\raggedright
\textbf{\textless{} 100ns}
\end{minipage} \\
\begin{minipage}[b]{\linewidth}\raggedright
\textbf{Corriente de fuga}
\end{minipage} & \begin{minipage}[b]{\linewidth}\raggedright
\textbf{\textless{} 10\ensuremath{\mu}A}
\end{minipage} \\
\midrule
\endhead
\bottomrule
\endlastfoot
\end{longtable}

\Needspace{6\baselineskip}
\captionof*{table}{\small\bfseries Tabla C-4.22. Especificaciones del RCD}

\begin{longtable}[]{@{}
  >{\raggedright\arraybackslash}p{(\columnwidth - 2\tabcolsep) * \real{0.3997}}
  >{\raggedright\arraybackslash}p{(\columnwidth - 2\tabcolsep) * \real{0.6003}}@{}}
\toprule
\begin{minipage}[b]{\linewidth}\raggedright
\textbf{Parámetro}
\end{minipage} & \begin{minipage}[b]{\linewidth}\raggedright
\textbf{Especificación}
\end{minipage} \\
\begin{minipage}[b]{\linewidth}\raggedright
\textbf{Corriente nominal}
\end{minipage} & \begin{minipage}[b]{\linewidth}\raggedright
\textbf{25A}
\end{minipage} \\
\begin{minipage}[b]{\linewidth}\raggedright
\textbf{Sensibilidad}
\end{minipage} & \begin{minipage}[b]{\linewidth}\raggedright
\textbf{30mA}
\end{minipage} \\
\begin{minipage}[b]{\linewidth}\raggedright
\textbf{Tipo}
\end{minipage} & \begin{minipage}[b]{\linewidth}\raggedright
\textbf{AC (corrientes alternas sinusoidales)}
\end{minipage} \\
\begin{minipage}[b]{\linewidth}\raggedright
\textbf{Tiempo de disparo}
\end{minipage} & \begin{minipage}[b]{\linewidth}\raggedright
\textbf{\ensuremath{\leq} 30ms a 5\ensuremath{\times}I\ensuremath{\Delta}n}
\end{minipage} \\
\begin{minipage}[b]{\linewidth}\raggedright
\textbf{Voltaje nominal}
\end{minipage} & \begin{minipage}[b]{\linewidth}\raggedright
\textbf{230V AC}
\end{minipage} \\
\begin{minipage}[b]{\linewidth}\raggedright
\textbf{Capacidad de ruptura}
\end{minipage} & \begin{minipage}[b]{\linewidth}\raggedright
\textbf{6kA}
\end{minipage} \\
\midrule
\endhead
\bottomrule
\endlastfoot
\end{longtable}

\Needspace{6\baselineskip}
\captionof*{table}{\small\bfseries Tabla C-5.1. Masa de Componentes (Kg)}

\begin{longtable}[]{@{}
  >{\raggedright\arraybackslash}p{(\columnwidth - 6\tabcolsep) * \real{0.4493}}
  >{\raggedright\arraybackslash}p{(\columnwidth - 6\tabcolsep) * \real{0.1963}}
  >{\raggedright\arraybackslash}p{(\columnwidth - 6\tabcolsep) * \real{0.1348}}
  >{\raggedright\arraybackslash}p{(\columnwidth - 6\tabcolsep) * \real{0.2196}}@{}}
\toprule
\begin{minipage}[b]{\linewidth}\raggedright
\textbf{COMPONENTE}
\end{minipage} & \begin{minipage}[b]{\linewidth}\raggedright
\textbf{MASA (kg)}
\end{minipage} & \begin{minipage}[b]{\linewidth}\raggedright
\textbf{\%TOTAL}
\end{minipage} & \begin{minipage}[b]{\linewidth}\raggedright
\textbf{UBICACIÓN}
\end{minipage} \\
\begin{minipage}[b]{\linewidth}\raggedright
\textbf{ESTRUCTURA PRINCIPAL}
\end{minipage} & \begin{minipage}[b]{\linewidth}\raggedright
\end{minipage} & \begin{minipage}[b]{\linewidth}\raggedright
\end{minipage} & \begin{minipage}[b]{\linewidth}\raggedright
\end{minipage} \\
\begin{minipage}[b]{\linewidth}\raggedright
Gabinete (Aluminio 5052-H32, 6mm)
\end{minipage} & \begin{minipage}[b]{\linewidth}\raggedright
21.976
\end{minipage} & \begin{minipage}[b]{\linewidth}\raggedright
46.4\%
\end{minipage} & \begin{minipage}[b]{\linewidth}\raggedright
Estructura externa
\end{minipage} \\
\begin{minipage}[b]{\linewidth}\raggedright
Rack y estructura interna (perfiles 30\ensuremath{\times}30mm)
\end{minipage} & \begin{minipage}[b]{\linewidth}\raggedright
5.700
\end{minipage} & \begin{minipage}[b]{\linewidth}\raggedright
12.0\%
\end{minipage} & \begin{minipage}[b]{\linewidth}\raggedright
Soporte interno
\end{minipage} \\
\begin{minipage}[b]{\linewidth}\raggedright
Aro/riel de giro (Aluminio 5052-H32)
\end{minipage} & \begin{minipage}[b]{\linewidth}\raggedright
7.800
\end{minipage} & \begin{minipage}[b]{\linewidth}\raggedright
16.5\%
\end{minipage} & \begin{minipage}[b]{\linewidth}\raggedright
Sistema visión base
\end{minipage} \\
\begin{minipage}[b]{\linewidth}\raggedright
Protector circular plano (Policarbonato, Ø420\ensuremath{\times}6mm)
\end{minipage} & \begin{minipage}[b]{\linewidth}\raggedright
1.000
\end{minipage} & \begin{minipage}[b]{\linewidth}\raggedright
2.1\%
\end{minipage} & \begin{minipage}[b]{\linewidth}\raggedright
Base sistema visión
\end{minipage} \\
\begin{minipage}[b]{\linewidth}\raggedright
\textbf{Subtotal estructura}
\end{minipage} & \begin{minipage}[b]{\linewidth}\raggedright
\textbf{36.476}
\end{minipage} & \begin{minipage}[b]{\linewidth}\raggedright
\textbf{77.0\%}
\end{minipage} & \begin{minipage}[b]{\linewidth}\raggedright
\end{minipage} \\
\begin{minipage}[b]{\linewidth}\raggedright
\end{minipage} & \begin{minipage}[b]{\linewidth}\raggedright
\end{minipage} & \begin{minipage}[b]{\linewidth}\raggedright
\end{minipage} & \begin{minipage}[b]{\linewidth}\raggedright
\end{minipage} \\
\begin{minipage}[b]{\linewidth}\raggedright
\textbf{SISTEMA DE MOVIMIENTO}
\end{minipage} & \begin{minipage}[b]{\linewidth}\raggedright
\end{minipage} & \begin{minipage}[b]{\linewidth}\raggedright
\end{minipage} & \begin{minipage}[b]{\linewidth}\raggedright
\end{minipage} \\
\begin{minipage}[b]{\linewidth}\raggedright
Motor PAP NEMA 23
\end{minipage} & \begin{minipage}[b]{\linewidth}\raggedright
1.500
\end{minipage} & \begin{minipage}[b]{\linewidth}\raggedright
3.2\%
\end{minipage} & \begin{minipage}[b]{\linewidth}\raggedright
Actuador principal
\end{minipage} \\
\begin{minipage}[b]{\linewidth}\raggedright
Biela (Aluminio 5052-H32)
\end{minipage} & \begin{minipage}[b]{\linewidth}\raggedright
0.142
\end{minipage} & \begin{minipage}[b]{\linewidth}\raggedright
0.3\%
\end{minipage} & \begin{minipage}[b]{\linewidth}\raggedright
Transmisión
\end{minipage} \\
\begin{minipage}[b]{\linewidth}\raggedright
Conector cámara-biela

(Aluminio 5052-H32)
\end{minipage} & \begin{minipage}[b]{\linewidth}\raggedright
0.857
\end{minipage} & \begin{minipage}[b]{\linewidth}\raggedright
1.8\%
\end{minipage} & \begin{minipage}[b]{\linewidth}\raggedright
Acoplamiento
\end{minipage} \\
\begin{minipage}[b]{\linewidth}\raggedright
Rodamiento rígido de bolas Ø120mm
\end{minipage} & \begin{minipage}[b]{\linewidth}\raggedright
1.000
\end{minipage} & \begin{minipage}[b]{\linewidth}\raggedright
2.1\%
\end{minipage} & \begin{minipage}[b]{\linewidth}\raggedright
Aro giratorio
\end{minipage} \\
\begin{minipage}[b]{\linewidth}\raggedright
\textbf{Subtotal sistema movimiento}
\end{minipage} & \begin{minipage}[b]{\linewidth}\raggedright
\textbf{3.499}
\end{minipage} & \begin{minipage}[b]{\linewidth}\raggedright
\textbf{7.4\%}
\end{minipage} & \begin{minipage}[b]{\linewidth}\raggedright
\end{minipage} \\
\begin{minipage}[b]{\linewidth}\raggedright
\end{minipage} & \begin{minipage}[b]{\linewidth}\raggedright
\end{minipage} & \begin{minipage}[b]{\linewidth}\raggedright
\end{minipage} & \begin{minipage}[b]{\linewidth}\raggedright
\end{minipage} \\
\begin{minipage}[b]{\linewidth}\raggedright
\textbf{COMPONENTES ELECTRÓNICOS POTENCIA}
\end{minipage} & \begin{minipage}[b]{\linewidth}\raggedright
\end{minipage} & \begin{minipage}[b]{\linewidth}\raggedright
\end{minipage} & \begin{minipage}[b]{\linewidth}\raggedright
\end{minipage} \\
\begin{minipage}[b]{\linewidth}\raggedright
UPS (500VA, baterías VRLA)
\end{minipage} & \begin{minipage}[b]{\linewidth}\raggedright
3.500
\end{minipage} & \begin{minipage}[b]{\linewidth}\raggedright
7.4\%
\end{minipage} & \begin{minipage}[b]{\linewidth}\raggedright
Nivel superior rack
\end{minipage} \\
\begin{minipage}[b]{\linewidth}\raggedright
SMPS S-360-24 (360W)
\end{minipage} & \begin{minipage}[b]{\linewidth}\raggedright
1.200
\end{minipage} & \begin{minipage}[b]{\linewidth}\raggedright
2.5\%
\end{minipage} & \begin{minipage}[b]{\linewidth}\raggedright
Nivel superior rack
\end{minipage} \\
\begin{minipage}[b]{\linewidth}\raggedright
Driver DM542T
\end{minipage} & \begin{minipage}[b]{\linewidth}\raggedright
0.350
\end{minipage} & \begin{minipage}[b]{\linewidth}\raggedright
0.7\%
\end{minipage} & \begin{minipage}[b]{\linewidth}\raggedright
Nivel inferior rack
\end{minipage} \\
\begin{minipage}[b]{\linewidth}\raggedright
\textbf{Subtotal electrónica potencia}
\end{minipage} & \begin{minipage}[b]{\linewidth}\raggedright
\textbf{5.050}
\end{minipage} & \begin{minipage}[b]{\linewidth}\raggedright
\textbf{10.7\%}
\end{minipage} & \begin{minipage}[b]{\linewidth}\raggedright
\end{minipage} \\
\begin{minipage}[b]{\linewidth}\raggedright
\end{minipage} & \begin{minipage}[b]{\linewidth}\raggedright
\end{minipage} & \begin{minipage}[b]{\linewidth}\raggedright
\end{minipage} & \begin{minipage}[b]{\linewidth}\raggedright
\end{minipage} \\
\begin{minipage}[b]{\linewidth}\raggedright
\textbf{COMPONENTES ELECTRÓNICOS CONTROL}
\end{minipage} & \begin{minipage}[b]{\linewidth}\raggedright
\end{minipage} & \begin{minipage}[b]{\linewidth}\raggedright
\end{minipage} & \begin{minipage}[b]{\linewidth}\raggedright
\end{minipage} \\
\begin{minipage}[b]{\linewidth}\raggedright
PCB Controlador (incluye base PCB)
\end{minipage} & \begin{minipage}[b]{\linewidth}\raggedright
0.050
\end{minipage} & \begin{minipage}[b]{\linewidth}\raggedright
0.1\%
\end{minipage} & \begin{minipage}[b]{\linewidth}\raggedright
Nivel inferior rack
\end{minipage} \\
\begin{minipage}[b]{\linewidth}\raggedright
Raspberry Pi 5 (montado en PCB)
\end{minipage} & \begin{minipage}[b]{\linewidth}\raggedright
0.075
\end{minipage} & \begin{minipage}[b]{\linewidth}\raggedright
0.2\%
\end{minipage} & \begin{minipage}[b]{\linewidth}\raggedright
Nivel inferior rack
\end{minipage} \\
\begin{minipage}[b]{\linewidth}\raggedright
HAT 4G LTE Quectel EG25-G (montado)
\end{minipage} & \begin{minipage}[b]{\linewidth}\raggedright
0.050
\end{minipage} & \begin{minipage}[b]{\linewidth}\raggedright
0.1\%
\end{minipage} & \begin{minipage}[b]{\linewidth}\raggedright
Nivel inferior rack
\end{minipage} \\
\begin{minipage}[b]{\linewidth}\raggedright
PCB de distribución de voltaje
\end{minipage} & \begin{minipage}[b]{\linewidth}\raggedright
0.300
\end{minipage} & \begin{minipage}[b]{\linewidth}\raggedright
0.6\%
\end{minipage} & \begin{minipage}[b]{\linewidth}\raggedright
Nivel inferior rack
\end{minipage} \\
\begin{minipage}[b]{\linewidth}\raggedright
\textbf{Subtotal electrónica control}
\end{minipage} & \begin{minipage}[b]{\linewidth}\raggedright
\textbf{0.475}
\end{minipage} & \begin{minipage}[b]{\linewidth}\raggedright
\textbf{1.0\%}
\end{minipage} & \begin{minipage}[b]{\linewidth}\raggedright
\end{minipage} \\
\begin{minipage}[b]{\linewidth}\raggedright
\end{minipage} & \begin{minipage}[b]{\linewidth}\raggedright
\end{minipage} & \begin{minipage}[b]{\linewidth}\raggedright
\end{minipage} & \begin{minipage}[b]{\linewidth}\raggedright
\end{minipage} \\
\begin{minipage}[b]{\linewidth}\raggedright
\textbf{PROTECCIÓN ELÉCTRICA}
\end{minipage} & \begin{minipage}[b]{\linewidth}\raggedright
\end{minipage} & \begin{minipage}[b]{\linewidth}\raggedright
\end{minipage} & \begin{minipage}[b]{\linewidth}\raggedright
\end{minipage} \\
\begin{minipage}[b]{\linewidth}\raggedright
Breaker C60N-C6
\end{minipage} & \begin{minipage}[b]{\linewidth}\raggedright
0.100
\end{minipage} & \begin{minipage}[b]{\linewidth}\raggedright
0.2\%
\end{minipage} & \begin{minipage}[b]{\linewidth}\raggedright
Nivel superior rack
\end{minipage} \\
\begin{minipage}[b]{\linewidth}\raggedright
RCD ID K 25A/30mA
\end{minipage} & \begin{minipage}[b]{\linewidth}\raggedright
0.150
\end{minipage} & \begin{minipage}[b]{\linewidth}\raggedright
0.3\%
\end{minipage} & \begin{minipage}[b]{\linewidth}\raggedright
Nivel superior rack
\end{minipage} \\
\begin{minipage}[b]{\linewidth}\raggedright
DPS VAL-MS 230
\end{minipage} & \begin{minipage}[b]{\linewidth}\raggedright
0.080
\end{minipage} & \begin{minipage}[b]{\linewidth}\raggedright
0.2\%
\end{minipage} & \begin{minipage}[b]{\linewidth}\raggedright
Nivel superior rack
\end{minipage} \\
\begin{minipage}[b]{\linewidth}\raggedright
Conversor DC-DC 5V (LM2596)
\end{minipage} & \begin{minipage}[b]{\linewidth}\raggedright
0.040
\end{minipage} & \begin{minipage}[b]{\linewidth}\raggedright
0.1\%
\end{minipage} & \begin{minipage}[b]{\linewidth}\raggedright
Nivel superior rack
\end{minipage} \\
\begin{minipage}[b]{\linewidth}\raggedright
Conversor DC-DC 12V (LM2596)
\end{minipage} & \begin{minipage}[b]{\linewidth}\raggedright
0.040
\end{minipage} & \begin{minipage}[b]{\linewidth}\raggedright
0.1\%
\end{minipage} & \begin{minipage}[b]{\linewidth}\raggedright
Nivel superior rack
\end{minipage} \\
\begin{minipage}[b]{\linewidth}\raggedright
\textbf{Subtotal protección}
\end{minipage} & \begin{minipage}[b]{\linewidth}\raggedright
\textbf{0.410}
\end{minipage} & \begin{minipage}[b]{\linewidth}\raggedright
\textbf{0.9\%}
\end{minipage} & \begin{minipage}[b]{\linewidth}\raggedright
\end{minipage} \\
\begin{minipage}[b]{\linewidth}\raggedright
\end{minipage} & \begin{minipage}[b]{\linewidth}\raggedright
\end{minipage} & \begin{minipage}[b]{\linewidth}\raggedright
\end{minipage} & \begin{minipage}[b]{\linewidth}\raggedright
\end{minipage} \\
\begin{minipage}[b]{\linewidth}\raggedright
\textbf{SENSORES}
\end{minipage} & \begin{minipage}[b]{\linewidth}\raggedright
\end{minipage} & \begin{minipage}[b]{\linewidth}\raggedright
\end{minipage} & \begin{minipage}[b]{\linewidth}\raggedright
\end{minipage} \\
\begin{minipage}[b]{\linewidth}\raggedright
WCS1600 (sensor corriente)
\end{minipage} & \begin{minipage}[b]{\linewidth}\raggedright
0.015
\end{minipage} & \begin{minipage}[b]{\linewidth}\raggedright
0.03\%
\end{minipage} & \begin{minipage}[b]{\linewidth}\raggedright
Rack
\end{minipage} \\
\begin{minipage}[b]{\linewidth}\raggedright
URM06 (sensor ultrasónico RS485)
\end{minipage} & \begin{minipage}[b]{\linewidth}\raggedright
0.025
\end{minipage} & \begin{minipage}[b]{\linewidth}\raggedright
0.05\%
\end{minipage} & \begin{minipage}[b]{\linewidth}\raggedright
Rack
\end{minipage} \\
\begin{minipage}[b]{\linewidth}\raggedright
DS18B20 (sensor temperatura con sonda)
\end{minipage} & \begin{minipage}[b]{\linewidth}\raggedright
0.020
\end{minipage} & \begin{minipage}[b]{\linewidth}\raggedright
0.04\%
\end{minipage} & \begin{minipage}[b]{\linewidth}\raggedright
Rack
\end{minipage} \\
\begin{minipage}[b]{\linewidth}\raggedright
\textbf{Subtotal sensores}
\end{minipage} & \begin{minipage}[b]{\linewidth}\raggedright
\textbf{0.060}
\end{minipage} & \begin{minipage}[b]{\linewidth}\raggedright
\textbf{0.13\%}
\end{minipage} & \begin{minipage}[b]{\linewidth}\raggedright
\end{minipage} \\
\begin{minipage}[b]{\linewidth}\raggedright
\end{minipage} & \begin{minipage}[b]{\linewidth}\raggedright
\end{minipage} & \begin{minipage}[b]{\linewidth}\raggedright
\end{minipage} & \begin{minipage}[b]{\linewidth}\raggedright
\end{minipage} \\
\begin{minipage}[b]{\linewidth}\raggedright
\textbf{SISTEMA DE VISIÓN}
\end{minipage} & \begin{minipage}[b]{\linewidth}\raggedright
\end{minipage} & \begin{minipage}[b]{\linewidth}\raggedright
\end{minipage} & \begin{minipage}[b]{\linewidth}\raggedright
\end{minipage} \\
\begin{minipage}[b]{\linewidth}\raggedright
Cámara domo
\end{minipage} & \begin{minipage}[b]{\linewidth}\raggedright
0.600
\end{minipage} & \begin{minipage}[b]{\linewidth}\raggedright
1.3\%
\end{minipage} & \begin{minipage}[b]{\linewidth}\raggedright
Montaje móvil
\end{minipage} \\
\begin{minipage}[b]{\linewidth}\raggedright
Soporte de cámara (metálico)
\end{minipage} & \begin{minipage}[b]{\linewidth}\raggedright
0.050
\end{minipage} & \begin{minipage}[b]{\linewidth}\raggedright
0.1\%
\end{minipage} & \begin{minipage}[b]{\linewidth}\raggedright
Acoplamiento
\end{minipage} \\
\begin{minipage}[b]{\linewidth}\raggedright
\textbf{Subtotal visión}
\end{minipage} & \begin{minipage}[b]{\linewidth}\raggedright
\textbf{0.650}
\end{minipage} & \begin{minipage}[b]{\linewidth}\raggedright
\textbf{1.4\%}
\end{minipage} & \begin{minipage}[b]{\linewidth}\raggedright
\end{minipage} \\
\begin{minipage}[b]{\linewidth}\raggedright
\end{minipage} & \begin{minipage}[b]{\linewidth}\raggedright
\end{minipage} & \begin{minipage}[b]{\linewidth}\raggedright
\end{minipage} & \begin{minipage}[b]{\linewidth}\raggedright
\end{minipage} \\
\begin{minipage}[b]{\linewidth}\raggedright
\textbf{ACCESORIOS Y FIJACIÓN}
\end{minipage} & \begin{minipage}[b]{\linewidth}\raggedright
\end{minipage} & \begin{minipage}[b]{\linewidth}\raggedright
\end{minipage} & \begin{minipage}[b]{\linewidth}\raggedright
\end{minipage} \\
\begin{minipage}[b]{\linewidth}\raggedright
Cableado (potencia, señal, datos)
\end{minipage} & \begin{minipage}[b]{\linewidth}\raggedright
0.400
\end{minipage} & \begin{minipage}[b]{\linewidth}\raggedright
0.8\%
\end{minipage} & \begin{minipage}[b]{\linewidth}\raggedright
Distribuido
\end{minipage} \\
\begin{minipage}[b]{\linewidth}\raggedright
Antenas y conectores externos
\end{minipage} & \begin{minipage}[b]{\linewidth}\raggedright
0.300
\end{minipage} & \begin{minipage}[b]{\linewidth}\raggedright
0.6\%
\end{minipage} & \begin{minipage}[b]{\linewidth}\raggedright
Externos
\end{minipage} \\
\begin{minipage}[b]{\linewidth}\raggedright
Ventilador 120\ensuremath{\times}120mm (24V, 12W)
\end{minipage} & \begin{minipage}[b]{\linewidth}\raggedright
0.200
\end{minipage} & \begin{minipage}[b]{\linewidth}\raggedright
0.4\%
\end{minipage} & \begin{minipage}[b]{\linewidth}\raggedright
Superior gabinete
\end{minipage} \\
\begin{minipage}[b]{\linewidth}\raggedright
Tornillería M24 (8 unid, conexión estructural)
\end{minipage} & \begin{minipage}[b]{\linewidth}\raggedright
1.440
\end{minipage} & \begin{minipage}[b]{\linewidth}\raggedright
3.0\%
\end{minipage} & \begin{minipage}[b]{\linewidth}\raggedright
Gabinete-soporte
\end{minipage} \\
\begin{minipage}[b]{\linewidth}\raggedright
Tornillería M12 (30 unid, montaje rack)
\end{minipage} & \begin{minipage}[b]{\linewidth}\raggedright
1.050
\end{minipage} & \begin{minipage}[b]{\linewidth}\raggedright
2.2\%
\end{minipage} & \begin{minipage}[b]{\linewidth}\raggedright
Rack/componentes
\end{minipage} \\
\begin{minipage}[b]{\linewidth}\raggedright
Tornillería M3 (8 unid, fijación PCBs)
\end{minipage} & \begin{minipage}[b]{\linewidth}\raggedright
0.016
\end{minipage} & \begin{minipage}[b]{\linewidth}\raggedright
0.03\%
\end{minipage} & \begin{minipage}[b]{\linewidth}\raggedright
PCBs
\end{minipage} \\
\begin{minipage}[b]{\linewidth}\raggedright
\textbf{Subtotal accesorios}
\end{minipage} & \begin{minipage}[b]{\linewidth}\raggedright
\textbf{3.406}
\end{minipage} & \begin{minipage}[b]{\linewidth}\raggedright
\textbf{7.2\%}
\end{minipage} & \begin{minipage}[b]{\linewidth}\raggedright
\end{minipage} \\
\begin{minipage}[b]{\linewidth}\raggedright
\end{minipage} & \begin{minipage}[b]{\linewidth}\raggedright
\end{minipage} & \begin{minipage}[b]{\linewidth}\raggedright
\end{minipage} & \begin{minipage}[b]{\linewidth}\raggedright
\end{minipage} \\
\begin{minipage}[b]{\linewidth}\raggedright
\textbf{MASA TOTAL DEL SISTEMA}
\end{minipage} & \begin{minipage}[b]{\linewidth}\raggedright
\textbf{47.126 kg}
\end{minipage} & \begin{minipage}[b]{\linewidth}\raggedright
\textbf{100.0\%}
\end{minipage} & \begin{minipage}[b]{\linewidth}\raggedright
\end{minipage} \\
\midrule
\endhead
\bottomrule
\endlastfoot
\end{longtable}

% ---------------------------------------------------------------------
\section*{ANEXO D: Presupuesto detallado y planos}
\addcontentsline{toc}{section}{ANEXO D: Presupuesto detallado y planos}\label{anexo:planos}

El presupuesto consolidado por rubros (CAPEX y OPEX) se presenta en el Capítulo~\ref{cap:analisis-economico}.
El desglose detallado por componente —fabricación mecánica, sistema de visión, electrónica de control,
protección eléctrica, sensores, instalación y prorrateo de servicios de nube— se conserva como hoja de
cálculo en el material digital de la tesis.

Los planos de fabricación se conservan como archivos CAD digitales y comprenden, como mínimo: (i) el
gabinete modular; (ii) el soporte tipo rack interno; (iii) el mecanismo de orientación de la cámara; (iv) el
poste y el soporte del controlador semafórico; y (v) la zapata de cimentación. Los esquemáticos eléctricos
de alimentación, distribución y de la tarjeta de control se incluyen igualmente como material digital y
corresponden a esquemas de diseño, sin verificación de reglas de manufactura (DRC) ni archivos Gerber.
